\documentclass[11pt,a4paper]{article}
\usepackage[utf8]{inputenc}
\usepackage[T1]{fontenc}
\usepackage[margin=2.5cm]{geometry}
\usepackage{fancyhdr}
\usepackage{titlesec}
\usepackage{tabularx}
\usepackage{booktabs}
\usepackage{xcolor}
\usepackage{enumitem}
\usepackage{hyperref}
\usepackage{amsmath}
\usepackage[expansion=false]{microtype}

\newcommand{\br}{\discretionary{}{}{}}

\pagestyle{fancy}
\fancyhf{}
\renewcommand{\headrulewidth}{0.4pt}
\fancyhead[L]{\small UTEQ -- ISR-401 Ingenier\'ia de Requerimientos}
\fancyhead[R]{\small Manuscrito borrador -- CarniCore -- Entrega 3 (2A)}
\fancyfoot[C]{\thepage}
\titleformat{\section}{\normalfont\Large\bfseries}{\thesection}{1em}{}
\titleformat{\subsection}{\normalfont\large\bfseries}{\thesubsection}{1em}{}

\hypersetup{colorlinks=true,linkcolor=black,urlcolor=blue,
  pdftitle={Manuscrito Borrador - Deteccion automatica de patrones de ambiguedad - CarniCore},
  pdfauthor={Equipo CarniCore - UTEQ}}

\title{\textbf{Detecci\'on autom\'atica de patrones de ambig\"uedad en requisitos funcionales:\\ un estudio exploratorio en el dominio de trazabilidad c\'arnica}}
\author{Castro Baja\~na, A.\,O. \and Crespo Espinoza, K.\,O. \and Gamarra Araujo, E.\,X. \and P\'erez Ruiz, C.\,A. \and Quintero Gende, E.\,J.}
\date{Borrador -- 02 de agosto de 2026 -- No enviado a ninguna revista}

\begin{document}
\maketitle

\begin{center}
\fbox{\begin{minipage}{0.9\linewidth}
\small\textbf{Estado de este documento.} Borrador de manuscrito, no enviado a ninguna revista. Reporta \'unicamente los resultados que existen como evidencia primaria real a la fecha: la salida de un detector autom\'atico de patrones aplicado sobre los 25 requisitos funcionales (RF) del sistema CarniCore. El estudio \textbf{no incluye} comparaci\'on con un panel de personas expertas: el equipo no cont\'o con acceso a evaluadores externos independientes durante esta entrega, por lo que esa validaci\'on se deja expl\'icitamente fuera de alcance (Secci\'on~7) en lugar de simularse o estimarse.
\end{minipage}}
\end{center}

\begin{abstract}
\textbf{Contexto.} La ambig\"uedad textual en los requisitos de software es una causa documentada de retrabajo y de malentendidos entre partes interesadas. \textbf{Objetivo.} Caracterizar la presencia de patrones l\'exico-sint\'acticos de ambig\"uedad -- cuantificadores vagos, conjunciones m\'ultiples y voz pasiva sin agente expl\'icito -- en el conjunto completo de requisitos funcionales de CarniCore, un sistema real en desarrollo para la gesti\'on de trazabilidad, pesaje inteligente e inventario de una distribuidora de productos c\'arnicos en Ecuador. \textbf{M\'etodo.} Se implement\'o un detector basado en expresiones regulares con tres categor\'ias de patr\'on y se ejecut\'o sobre los 25 RF verbatim del ERS/SRS v3.0 (poblaci\'on completa, sin muestreo). \textbf{Resultados.} El detector no activ\'o ninguna de las tres categor\'ias en ninguno de los 25 RF (0/25, 0\,\%). \textbf{Conclusiones.} El hallazgo es consistente con un estilo de redacci\'on disciplinado y homog\'eneo en el corpus analizado, pero no permite concluir ausencia de ambig\"uedad en sentido amplio: los patrones detectados son \'unicamente l\'exico-sint\'acticos y no capturan la ambig\"uedad sem\'antica o pragm\'atica (p.\,ej., umbrales temporales no especificados), cuya identificaci\'on requiere juicio humano y queda como trabajo futuro no ejecutado en este estudio por falta de acceso a un panel de personas evaluadoras.
\end{abstract}

\section{Introducci\'on}

La calidad de los requisitos de software condiciona directamente el \'exito de un proyecto: la ambig\"uedad textual en un requisito -- la posibilidad de que admita m\'as de una interpretaci\'on razonable -- es una de las causas m\'as documentadas de retrabajo, defectos y malentendidos entre partes interesadas en Ingenier\'ia de Requisitos (Femmer et al., 2017; Ferrari y Esuli, 2019). Detectar esta ambig\"uedad de forma temprana, idealmente de manera automatizada, es un objetivo pr\'actico de larga data en el \'area.

Este art\'iculo presenta un estudio exploratorio que aplica un detector autom\'atico de ambig\"uedad basado en patrones textuales (cuantificadores vagos, conjunciones m\'ultiples y voz pasiva sin agente expl\'icito) sobre el conjunto completo de 25 requisitos funcionales de CarniCore, un sistema real en desarrollo. A diferencia de un dise\~no comparativo cl\'asico de detecci\'on de ambig\"uedad, este estudio no contrasta la salida del detector contra el juicio de personas expertas: esa comparaci\'on qued\'o fuera de alcance por falta de acceso a un panel de evaluadores independientes durante el periodo de esta entrega, y se documenta como trabajo futuro (Secci\'on~7).

\subsection*{Contribuciones}
\begin{itemize}[leftmargin=*]
\item Un detector de ambig\"uedad de c\'odigo abierto, reproducible, con tres categor\'ias de patr\'on documentadas (\texttt{detector\_ambiguedad.py}).
\item Una caracterizaci\'on emp\'irica real -- no simulada -- de la presencia de estos patrones en el conjunto completo de RF de un sistema en producci\'on temprana.
\item Un conjunto de datos abierto (Zenodo, licencia CC~BY~4.0) con la salida del detector sobre los 25 RF.
\item Una discusi\'on expl\'icita de las limitaciones de un detector puramente l\'exico frente a la ambig\"uedad sem\'antica, que motiva un dise\~no de validaci\'on con panel experto para trabajo futuro.
\end{itemize}

\section{Trabajo relacionado}

La detecci\'on autom\'atica de ambig\"uedad y de ``malos olores'' (\emph{smells}) en requisitos en lenguaje natural ha sido estudiada mediante enfoques basados en reglas y expresiones regulares (Femmer et al., 2017), en procesamiento de lenguaje natural cl\'asico (Ferrari y Esuli, 2019; Chantree et al., 2006) y, m\'as recientemente, en modelos de lenguaje de gran escala aplicados en contextos industriales (Bashir et al., 2025). Gervasi et al. (2019) proponen un marco unificador para distinguir tipos de ambig\"uedad (l\'exica, sint\'actica, sem\'antica, pragm\'atica), marco que se retoma en la Secci\'on~6 para interpretar los l\'imites de este estudio. Philippo et al. (2013) reportan que la ambig\"uedad percibida por personas evaluadoras no siempre coincide con la ambig\"uedad detectada formalmente -- un hallazgo directamente relevante para justificar por qu\'e el detector aqu\'i usado, al ser puramente l\'exico, no agota la noci\'on de ambig\"uedad de un requisito.

A diferencia de estudios que contrastan un detector contra el consenso de un panel experto (dise\~no que este equipo pre-registr\'o en OSF para una fase posterior, \url{https://osf.io/wud69}), este estudio se limita a la caracterizaci\'on descriptiva del detector sobre un conjunto de requisitos originales, elicitados de campo por el propio equipo para un sistema real -- lo que introduce validez externa distinta (dominio no explorado previamente en la literatura de ambig\"uedad de requisitos: trazabilidad de productos c\'arnicos en un contexto latinoamericano) a costa de no contar con una referencia de juicio humano en esta entrega.

\section{Dise\~no del estudio}

\textbf{Pregunta de investigaci\'on.} \textquestiondown En qu\'e medida los requisitos funcionales de CarniCore presentan los patrones l\'exico-sint\'acticos de ambig\"uedad descritos en la literatura de \emph{requirements smells} (cuantificadores vagos, conjunciones m\'ultiples, voz pasiva sin agente expl\'icito)?

\textbf{Dise\~no.} Estudio observacional, descriptivo, transversal y censal: se analiza el 100\,\% de la poblaci\'on de 25 RF de CarniCore (RF-01 a RF-25), sin muestreo. No se plantea contraste de hip\'otesis inferencial en esta entrega, dado que no hay una segunda fuente de clasificaci\'on (panel experto) contra la cual comparar al detector; el dise\~no comparativo completo permanece pre-registrado en OSF para ejecuci\'on futura.

\section{M\'etodo}

\subsection{Poblaci\'on}
Los 25 RF completos del ERS/SRS v3.0 de CarniCore, citados de forma verbatim (Secci\'on 3.2 del ERS/SRS), sin los atributos de identificador, fuente, prioridad ni estabilidad.

\subsection{Detector autom\'atico}
Script propio en Python (\texttt{06\_\br Experimento/\br scripts\_\br analisis/\br detector\_\br ambiguedad.py}), con tres categor\'ias de patr\'on:
\begin{itemize}[leftmargin=*]
\item \textbf{Cuantificadores vagos}: adjetivos y adverbios sin criterio verificable (p.\,ej. ``apropiado'', ``suficiente'', ``r\'apidamente'', ``aproximadamente'').
\item \textbf{Conjunciones m\'ultiples}: m\'as de tres conectores coordinantes (\emph{y}/\emph{o}) o m\'as de dos verbos modales (\emph{deber\'a}) dentro de un mismo RF, como se\~nal de un requisito potencialmente compuesto.
\item \textbf{Voz pasiva sin agente expl\'icito}: construcciones pasivas perifr\'asticas o reflejas no seguidas de un complemento agente (\emph{``por + sujeto''}).
\end{itemize}
Un RF se clasifica como \texttt{ambiguo\_detector = 1} si activa al menos una categor\'ia. La regla de decisi\'on y el listado completo de patrones se fijaron antes de examinar los resultados y no se modificaron despu\'es.

\section{Resultados}

Se ejecut\'o el detector sobre los 25 RF. Ning\'un RF activ\'o ninguna categor\'ia:

\begin{center}
\begin{tabular}{@{}lc@{}}
\toprule
Categor\'ia de patr\'on & RF que la activaron \\
\midrule
Cuantificador vago & 0/25 \\
Conjunci\'on m\'ultiple ($>$3 conectores \'o $>$2 verbos modales) & 0/25 \\
Voz pasiva sin agente expl\'icito & 0/25 \\
\midrule
\textbf{Al menos una categor\'ia (ambiguo\_detector = 1)} & \textbf{0/25 (0\,\%)} \\
\bottomrule
\end{tabular}
\end{center}

Este resultado es consistente con el estilo de redacci\'on observado en el ERS: los 25 RF siguen sistem\'aticamente la plantilla \emph{``El sistema deber\'a permitir/registrar/calcular + objeto directo concreto''}, en voz activa, sin los cuantificadores l\'exicos ni las estructuras pasivas que el detector busca. Es un resultado determin\'istico y reproducible: cualquier reejecuci\'on del script sobre el mismo texto produce id\'entica salida.

\section{Discusi\'on}

Un resultado de 0/25 en un detector puramente l\'exico \textbf{no equivale a afirmar que los 25 RF est\'an libres de ambig\"uedad}. Siguiendo el marco de Gervasi et al. (2019), el detector aqu\'i implementado cubre \'unicamente ambig\"uedad l\'exica y sint\'actica superficial; no cubre ambig\"uedad sem\'antica ni pragm\'atica. Dos ejemplos del propio corpus ilustran el l\'imite: RF-08 (``genera una alerta visual cuando \emph{est\'e pr\'oximo a vencer}'') no especifica un umbral num\'erico expl\'icito en su redacci\'on principal -- aunque el criterio de aceptaci\'on asociado s\'i lo fija en un documento separado del ERS --, y RF-21 (``sincroniz\'andolas \emph{autom\'aticamente} al restablecerse la conectividad'') no describe el mecanismo de sincronizaci\'on en la propia oraci\'on del requisito. Ninguno de los dos activa un patr\'on l\'exico del detector, pero ambos podr\'ian generar interpretaciones distintas en un lector sin acceso al resto del documento -- precisamente el tipo de discrepancia que Philippo et al. (2013) documentan entre ambig\"uedad formalmente detectada y ambig\"uedad percibida.

Esta discusi\'on es preliminar y no puede resolverse sin la fase comparativa contra juicio humano, que este estudio no ejecut\'o. El equipo considera que el resultado 0/25 es en s\'i mismo informativo -- sugiere que, si existe ambig\"uedad percibida en este corpus, es probable que sea de tipo sem\'antico o pragm\'atico y no l\'exico -- pero se abstiene de cuantificar esa probabilidad sin datos de un panel experto.

\section{Amenazas a la validez}
\begin{itemize}[leftmargin=*]
\item \textbf{Validez de constructo.} El detector solo opera sobre tres categor\'ias l\'exico-sint\'acticas; un resultado de 0 positivos puede reflejar aus\'encia real de esos patrones o cobertura insuficiente del detector (p.\,ej., listas de cuantificadores vagos incompletas para el espa\~nol t\'ecnico usado en este dominio).
\item \textbf{Ausencia de validaci\'on externa.} No se compar\'o la salida del detector contra ning\'un juicio humano en esta entrega; no hay, por tanto, evidencia de precisi\'on o \emph{recall} del detector frente a la ambig\"uedad realmente percibida por lectores del ERS.
\item \textbf{Tama\~no y alcance.} 25 RF de un \'unico sistema, elicitados por un \'unico equipo, en un \'unico dominio (trazabilidad c\'arnica) y en espa\~nol de Ecuador; no se reclama generalizaci\'on a otros corpus o idiomas.
\item \textbf{Ausencia de revisi\'on por pares.} Este es un borrador no enviado ni revisado.
\end{itemize}

\section{Trabajo futuro}

El equipo pre-registr\'o en OSF (\url{https://osf.io/wud69}, 02/08/2026) un dise\~no comparativo completo -- detector vs.\ consenso de un m\'inimo de tres personas expertas, con c\'alculo de precisi\'on, \emph{recall}, F1 y acuerdo inter-evaluador ($\kappa$ de Cohen y de Fleiss) -- que no se ejecut\'o en el periodo de esta entrega por falta de acceso a un panel de evaluadores independientes. Ese dise\~no permanece disponible en el protocolo experimental (\texttt{06\_\br Experimento/\br protocolo.pdf}) para ejecuci\'on en una fase posterior del proyecto, en caso de contar con los recursos humanos necesarios.

\section*{Disponibilidad de datos}
El conjunto de datos (salida real del detector sobre los 25 RF) se deposita en Zenodo bajo licencia Creative Commons Atribuci\'on 4.0 Internacional (CC~BY~4.0), conforme a \texttt{07\_\br Publicacion/\br dataset\_\br zenodo/}. El DOI se agrega tras el dep\'osito.

\section*{Referencias (parciales, formato a normalizar antes del env\'io)}
{\small
\begin{itemize}[leftmargin=*]
\item Bashir, S., Ferrari, A., Khan, M.\,A., Strandberg, P.\,E., Haider, Z., Saadatmand, M., Bohlin, M. (2025). Requirements Ambiguity Detection and Explanation with LLMs: An Industrial Study. \emph{ICSME 2025}.
\item Chantree, F., Nuseibeh, B., De Roeck, A., Willis, A. (2006). Identifying Nocuous Ambiguities in Natural Language Requirements.
\item Femmer, H., M\'endez Fern\'andez, D., Wagner, S., Eder, S. (2017). Rapid quality assurance with requirements smells. \emph{Journal of Systems and Software}, 123, 190--213.
\item Ferrari, A., Esuli, A. (2019). An NLP approach for cross-domain ambiguity detection in requirements engineering. \emph{Automated Software Engineering}, 26(3), 559--598.
\item Gervasi, V., Ferrari, A., Zowghi, D., Spoletini, P. (2019). Ambiguity in Requirements Engineering: Towards a Unifying Framework. \emph{From Software Engineering to Formal Methods and Tools, and Back}. Springer.
\item Philippo, E.\,J., Heijstek, W., Kruiswijk, B., Chaudron, M.\,R.\,V., Berry, D.\,M. (2013). Requirement Ambiguity Not as Important as Expected -- Results of an Empirical Evaluation. \emph{REFSQ Workshops}, 65--79.
\end{itemize}
}

\vspace{0.8cm}
\noindent\rule{\linewidth}{0.4pt}
\small
\textbf{Nota de cumplimiento.} Todos los datos reportados en este manuscrito (Secci\'on 5) son reales, producidos por ejecuci\'on efectiva de \texttt{detector\_ambiguedad.py} sobre los 25 RF verbatim del ERS/SRS. No se reporta ninguna cifra de precisi\'on, \emph{recall}, F1 ni $\kappa$ porque la fase comparativa con panel experto no se ejecut\'o en este periodo -- se documenta como trabajo futuro (Secci\'on~7) en lugar de estimarse o simularse, conforme a la advertencia de la Gu\'ia y R\'ubrica de la Entrega 3 (2A) sobre fabricaci\'on de evidencia (gatekeeper G4).

\end{document}
